% -*- coding: utf-8 -*-
% Preámbulo común para la asignatura Lógica (Grado en Filosofía, Universidad de Sevilla)
% Compatible tanto para compilación independiente (article) como en el manual conjunto (book/report)

\usepackage[T1]{fontenc}
\usepackage[utf8]{inputenc}
\usepackage[spanish,es-tabla]{babel}
\usepackage{amsmath,amssymb,amstext,amsthm,array}
\usepackage[top=20mm,bottom=20mm,left=20mm,right=20mm,headsep=5mm,footskip=10mm]{geometry}
\usepackage{enumerate}
\usepackage{tikz}
\usetikzlibrary{trees}

% Paquete Fitch y reglas canónicas de la asignatura
% Altura de línea algo menor que la predeterminada (4.5ex/3.5ex/1.5ex) para que
% las derivaciones sean más compactas y den menos problemas de paginación.
\usepackage[justsep=1em,rules=myrules,justformat=myjust,height=3.8ex,topheight=3.2ex,depth=1.3ex]{fitch}
\newcommand{\myjust}[2]{#1 #2}
\newcommand{\myrules}{
 \ndrules
 \def\ci{\by{$\to$I}}
 \def\ce{\by{$\to$E}}
 \def\x{\by{X}}
 \def\ip{\by{IP}}
 \def\ds{\by{DS}}
 \def\mt{\by{MT}}
 \def\lem{\by{LEM}}
 \def\dne{\by{DNE}}
 \def\dem{\by{DeM}}
 \def\cq{\by{CQ}}
 \def\bi{\by{$\leftrightarrow$I}}
 \def\be{\by{$\leftrightarrow$E}}
 \def\idi{\by{$=$I}}
 \def\ide{\by{$=$E}}}

% Enlaces hipertexto
\usepackage[hidelinks,colorlinks=true,linkcolor=blue!60!black,urlcolor=blue!60!black]{hyperref}

% Aplicaciones web de apoyo (GitHub Pages del repositorio fersoler/AplicacionesLogica).
% Si cambia la dirección, basta cambiarla aquí.
\newcommand{\urlaplicaciones}{https://fersoler.github.io/AplicacionesLogica/}
\newcommand{\webaplicaciones}{\expandafter\url\expandafter{\urlaplicaciones}}
\newcommand{\enlaceapp}[2]{\href{\urlaplicaciones #1/}{\emph{#2}}}
\newcommand{\appsintaxis}{\enlaceapp{sintaxis}{Sintaxis de las fórmulas}}
\newcommand{\appdeduccion}{\enlaceapp{deduccion}{Deducción natural}}
\newcommand{\apptablas}{\enlaceapp{tablas}{Tablas de verdad}}
\newcommand{\appkalmar}{\enlaceapp{kalmar}{La prueba de Kalmár}}
\newcommand{\appmundo}{\enlaceapp{mundo}{El mundo de Tarski}}
\newcommand{\apptableaux}{\enlaceapp{tableaux}{Tableaux semánticos}}
% Nota breve tras una explicación: con qué aplicación se practica. Empieza con
% el icono de un portátil (fontawesome5) en el color de los enlaces.
\usepackage{fontawesome5}
\newcommand{\practica}[1]{\par\smallskip\noindent{\small{\color{blue!60!black}\faLaptop}\enspace\textit{Para practicar.} #1\par}\smallskip}

% Teoremas y entornos formales
\ifcsname chapter\endcsname
  \newtheorem{theorem}{Teorema}[chapter]
\else
  \newtheorem{theorem}{Teorema}[section]
\fi
\newtheorem{lemma}[theorem]{Lema}
\newtheorem{definition}[theorem]{Definición}
\newtheorem{definicion}[theorem]{Definición}
\newtheorem{example}[theorem]{Ejemplo}
\newtheorem{ejemplo}[theorem]{Ejemplo}
\newtheorem{corollary}[theorem]{Corolario}
\newtheorem{proposition}[theorem]{Proposición}
\newtheorem{observacion}[theorem]{Observación}
\newtheorem{nota}[theorem]{Nota}
\newtheorem{aplicacion}[theorem]{Aplicación}

\newenvironment{resumenbox}{%
  \begin{quote}\small
}{%
  \end{quote}
}

% Ejercicios: en el manual (clase book) se numeran por tema (Ejercicio 4.12).
% El título entre paréntesis va en redonda, como en los teoremas y ejemplos.
\newcounter{preg}
\ifcsname chapter\endcsname
  \makeatletter\@addtoreset{preg}{chapter}\makeatother
  \renewcommand{\thepreg}{\thechapter.\arabic{preg}}
\fi
\newenvironment{pregunta}[1][]{%
  \par\addvspace{0.4cm}\refstepcounter{preg}\noindent\textbf{Ejercicio \thepreg}\if\relax\detokenize{#1}\relax\textbf{.}\else\ (#1)\textbf{.}\fi\enspace\ignorespaces
}{%
  \par\vspace{0.4cm}%
}

% Soluciones (apéndice de soluciones)
% Dentro de una solución no se permite un salto de página entre el título y lo
% que le sigue (por ejemplo, una derivación en un entorno center), ni que un
% párrafo deje una sola línea al pie o al principio de una página.
\newenvironment{solucion}[1]{%
  \par\medskip
  \csname @beginparpenalty\endcsname=10000\relax
  \clubpenalty=10000\relax\widowpenalty=10000\relax
  \noindent\textbf{Solución del ejercicio~\ref{#1}.}\enspace\ignorespaces
}{%
  \par\medskip%
}

% Derivación dentro de un recuadro de regla (tabla con \hline): deja un poco
% de aire para que la línea vertical de la derivación no toque las líneas
% horizontales de la tabla.
\newcommand{\cajaregla}[1]{\vbox{\vskip 5pt\hbox{#1}\vskip 4pt}}

% Marca para secciones de ampliación (no necesarias para seguir el curso)
\newcommand{\ampliacion}{\texorpdfstring{\,$^{\star}$}{ (*)}}

\newcounter{pregexamen}
\newenvironment{preguntaexamen}[2]{%
  \refstepcounter{pregexamen}%
  \vspace{0.4cm}%
  \noindent\textbf{Pregunta \arabic{pregexamen}} \hfill \textbf{[#1 puntos]}\\[0.15cm]%
  \noindent\textit{#2}\\[0.2cm]%
}{%
  \vspace{0.3cm}%
}

\newtheorem{ejercicio}{Ejercicio}

\newcounter{ejerciciocnt}
\newcommand{\ejer}[1]{\vspace{1.5em}\noindent\stepcounter{ejerciciocnt}\textbf{Ejercicio \theejerciciocnt.} #1\vspace{0.5em}}

\newenvironment{resumentema}{\begin{quote}\small\textbf{Resumen.}\enspace}{\end{quote}\bigskip}
